\section{Contributions}

The main contribution is a reproducible protocol for converting a legal source
case into a clustered, rubric-scored evaluation of LLM legal reasoning. The
protocol differs from static legal benchmarks in three ways.

\textbf{Source-grounded benchmark synthesis.} The benchmark instance is
generated from a citable source case. The source-derived packet contains the
question, gold answer, controlling legal gates, boundary perturbations, and
provenance hashes used downstream.

\textbf{Case-specific rubric induction.} The scoring rubric is generated from
the locked source packet rather than reused from a generic answer-quality
checklist. Each row must correspond to a legally observable component of the
case theory and include source support.

\textbf{Reasoning-centroid evaluation.} The pipeline clusters natural model
answers by normalized legal-reasoning signatures before judging. A
rubric-conditioned judge scores cluster representatives, and centroid scores are
projected to all members to produce model-level rankings with row-level
rationales.

These contributions separate three operations that are frequently conflated in
LLM evaluation: constructing a legal benchmark, identifying families of legal
reasoning in model outputs, and scoring those families under a source-grounded
rubric.

\subsection{Statute-of-Frauds Litmus Test}

Statute of Frauds is the first litmus domain because it naturally contains
discrete legal gates and boundary facts: marriage consideration, one-year
performance, suretyship, land interests, goods/UCC writings, executor or
administrator promises, writing substitutes, and equitable exceptions. That
structure makes it a useful early test for whether an automated pipeline can
generate legally meaningful variations and separate different reasoning paths.

In the present validation, the strongest evidence concerns pipeline mechanics
and representative Statute-of-Frauds reasoning paths. The live model batch
exercises the full source-to-ranking path on a marriage-consideration
beneficiary-certificate case, but because the sampled models converged on one
wife/certificate theory, it is not a sufficient test of Dasha's
divergence-detection claim. The 500-response regression suite checks the Dasha
clustering container against eight Statute-of-Frauds archetypes:
wife/certificate enforcement,
association replacement-certificate control, one-year/no-writing bar,
one-year-under-year enforceability, suretyship/main-purpose exception, land
part-performance, goods/UCC quantity defect, and executor personal-promise
without writing. This does not prove every possible Statute-of-Frauds case is
solved, but it shows that the pipeline's schemas, clustering machinery, and
validation surfaces cover the major gate families needed for that domain.

\subsection{Generalization Beyond Statute of Frauds}

The architecture is doctrine-general because every live stage is driven by the
source case and canonical instruction context. Frank is asked to infer the
doctrine and legal gates; Karthic builds a rubric from the locked Frank packet;
Dasha normalizes whatever doctrine, trigger, outcome, and exception structure is
present in the responses; Judge scores against Karthic's rows rather than a
Statute-of-Frauds-specific checklist. To apply the method to another legal
field, the expected change is not the pipeline architecture but the instruction
context and calibration examples supplied to the agents. Broader legal-domain
claims therefore require additional source packs and held-out cases, but the
pipeline is already structured so those additions are data and instruction
extensions rather than a new evaluation system.
